\section{Review Method and Scope}\label{sec:review-method-and-scope}

This chapter reviews the literature needed to support the research questions
posed in Section~\ref{sec:research-questions}. The review is deliberately
targeted rather than bibliometric: its purpose is to establish (1) how
language-model copilots for security operations behave and how they validate
their own output; (2) the attack surface of applications that integrate
language models with tooling; (3) the available methods for detecting model
fabrication and for measuring answer agreement or uncertainty; and (4) which of
these techniques remain viable under the deployment constraints adopted in this
dissertation---an offline, CPU-only, single-inference budget with no retrieval
at query time. These four questions map directly onto RQ1--RQ3.

The corpus is drawn from the computing literature and from authoritative
technical sources. Peer-reviewed and established-venue material was sought
across the ACM proceedings (including the ESEC/FSE joint conference and the
Artificial Intelligence and Security (AISec) workshop), the empirical NLP and
machine learning conferences (EMNLP, ICLR, NeurIPS), the ACM Computing Surveys
journal, and the arXiv preprint archive; standards and tooling references were taken
from their canonical sources (the IETF RFC series, the OWASP project, and the
official documentation of the Nmap and Metasploit tools). Candidate works were
located with controlled query themes covering the joint areas of interest:
security copilots and penetration testing assisted by language models, prompt
injection and indirect prompt injection, hallucination detection, semantic
uncertainty, self-consistent decoding, retrieval-augmented generation, and
authorisation control for tool-driven systems.

Works were included when they satisfied two criteria: relevance to at least one
research question, and a verifiable primary source reporting either an empirical
evaluation or an authoritative specification. The machine-learning literature
covered is concentrated in the period 2019--2025, which spans the advent of the
sentence-embedding and language-model techniques relied on here
\parencite{reimers2019sbert,lewis2020rag}; older technical standards and
tooling references are included regardless of year because they define stable
interfaces rather than evolving research claims. Preprint versions were
preferred only when a peer-reviewed version was not available for a work; each
entry in the reference list was verified against its real source before it was
cited in this dissertation. The review is bounded by the deployment constraint:
methods that assume interactive generation, a cloud-hosted verifier, or
retrieval infrastructure are surveyed because they are the competing
alternatives against which the guarded pipeline must be justified, and their
unsuitability for the offline single-inference budget is part of the analysis
in Section~\ref{sec:gap-synthesis}.

The remainder of the chapter is organised as follows.
Section~\ref{sec:domain-foundations} establishes the theoretical background of
intent recognition and of the guarantees expected from a safety-critical
interface. Section~\ref{sec:related-research} examines the related research on
LLM copilots, attacks, and fabrication detection in detail.
Section~\ref{sec:existing-systems} reviews existing systems, products and
techniques against the same criteria. Section~\ref{sec:gap-synthesis} synthesises
the findings into the gap that motivates this dissertation and into the
conceptual framework that the design chapters build on.

\TODO{Confirm whether the faculty's assessment rubric requires a formal
systematic-review protocol (screening counts, PRISMA-style reporting). The
current review is a targeted narrative review; screening statistics are not
reported because no formal screening process was conducted.}